\documentclass[11pt]{article}

\def\chapitre{17}
\def\pagetitle{Variables aléatoires.}

\input{setup.tex}

\begin{document}

\input{title.tex}

\vspace*{-0.7cm}

\section{Généralités.}

\begin{defi}{Variable aléatoire discrète.}{}
    Soit $(\O,\A,\Pr)$ un espace probabilisé et $X:\O\to\R$ une application.\\
    On dit que $X$ est une \bf{variable aléatoire discrète} (v.a.d.) sur $(\O,\A,\Pr)$ ssi :
    \begin{enumerate}
        \item $X(\O)$ est au plus dénombrable.
        \item $\forall x \in X(\O), ~ X^{-1}(\{x\})\in\A$, noté $(X = x)$.
    \end{enumerate}
\end{defi}

\vspace*{-0.8cm}

\begin{ex}{}{}
    On tire un dé à 6 faces deux fois. On note $X$ le plus grand numéro obtenu.\\
    Si $\w=(1, 5)$, alors $X(\w)=5$; si $\w=(3,2)$, alors $X(\w)=3$.
\end{ex}

\vspace*{-0.8cm}

\begin{prop}{$(X\in A)$}{}
    Soit $X$ une v.a.d. sur $(\O,\A,\Pr)$. Alors $\forall A \subset X(\O), ~ (X \in A)\in\A$.
    \tcblower
    Soit $A\in X(\O)$. On a $(X\in A) \iff \exists a \in A, ~ X=a$ donc $(X\in A) = \bigcup_{a\in A}(X=a)\in \A$.\\
    En effet, c'est un réunion au plus dénombrable d'évènements de $\A$.\\
    En outre, par $\s$-additivité (les évènements sont 2-à-2 incompatibles) :
    \begin{equation*}
        \Pr(X\in A) = \sum_{a\in A}\Pr(X=a).
    \end{equation*}
\end{prop}

\vspace*{-0.8cm}

\begin{ex}{Cas particuliers.}{}
    \begin{enumerate}
        \item $A=[a,+\infty[~\subset\R$ : $(X\in A) = (X\geq a) = \bigcup_{\substack{x\geq a\\x\in X(\O)}}(X=x)$.
        \item $A=]-\infty,a]~\subset\R$ : $(X \in A) = (X < a) = \bigcup_{\substack{x < a \\ x \in X(\O)}}(X=x)$. 
    \end{enumerate}
    Supposons de plus que $X(\O) = \N$, alors $\forall n \in \N, ~ (X\geq n) = \bigcup_{k\geq n}(X=k)=(X=n)\sqcup(X\geq n+1)$.\\
    Par $\s$-additivité, $\Pr(X\geq n) = \Pr(X=n)+\Pr(X\geq n+1)$.
\end{ex}

\pagebreak

\begin{ex}{}{}
    Soit $(\O,\A,\Pr)$ un espace probabilisé et $A\in\A$. Vérifier que $\1_A$ est une v.a.d. sur $(\O,\A,\Pr)$.
    \tcblower
    On a $\1_A(\O)\subset\{0,1\}$ donc au plus dénombrable car fini.\\
    Puis $(\1_A = 0) = \ov{A} \in \A$ et $(\1_A = 1) = A \in \A$ donc $\forall x \in \1_A(\O), ~ (\1_A = x) \in \A$. C'est une v.a.d. 
\end{ex}

\begin{exercice}{}{}
    Soit $(\O,\A,\Pr)$ un espace probabilisé et $N$ une v.a.d. sur cet espace à valeurs dans $\N$.\\
    On prend $(X_n)_{n\in\N}$ une suite de v.a.d. sur $(\O,\A,\Pr)$, et $Y=X_N$. Montrer que $Y$ est une v.a.d.
    \tcblower
    Déjà, $Y:\O\to\R$ avec $Y(\w)=X_{N(\w)}(\w)$, donc $Y(\O)\subset\bigcup_{n\in\N}X_n(\O)$.\\
    Soit $y\in Y(\O)$, alors $(Y = y) = \bigcup_{n\in\N}(N = n \cap X_n = y)\in\A$.
\end{exercice}

\begin{defi}{Loi d'une variable aléatoire.}{}
    Soit $(\O,\A,\P)$ un espace probabilisé et $X$ une v.a.d. sur cet espace.\\
    On pose $\forall A \in X(\O), ~ \Pr_X(A)=\Pr(X\in A)$ la loi de $X$, c'est une probabilité sur $(X(\O), \P(X(\O)))$.
    \tcblower
    Montrons que $\Pr_X$ est bien une probabilité sur $(X(\O), \P(X(\O)))$.\\
    $\bullet$ $\Pr_X(X(\O))=\Pr(X\in X(\O))=\Pr(\O)=1$.\\
    $\bullet$ $\forall A \in \P(X(\O))$, $\Pr_X(A)=\Pr(X\in A)\geq 0$.\\
    $\bullet$ Pour $(A_n)_n\in\P(X(\O))^\N$ 2-à-2 incompatibles.
    \begin{equation*}
        \Pr_X(\bigcup_{n\in\N}A_n) = \Pr(X\in\bigcup_{n\in\N}A_n) = \Pr(\bigcup_{n\in\N}(X\in A_n)) = \sum_{n\in\N}\Pr(X\in A_n).
    \end{equation*}
\end{defi}

\begin{defi}{Mêmes lois.}{}
    Soient $X,Y$ deux v.a.d. sur $(\O,\A,\Pr)$. On dit que $X$ et $Y$ ont même loi, et on note $X\sim Y$ ssi $\Pr_X = \Pr_Y$.\\
    Alors $X$ et $Y$ ont même probabilité sur tout évènement.
\end{defi}

\begin{center}{}{}
    \fbox{$X\sim Y \not\ra \forall \w \in \O, ~ X(\w) = Y(\w)$}
\end{center}

\pagebreak

\begin{prop}{}{}
    Soit $(\O, A, \Pr)$ un espace probabilisé et $X,Y$ deux v.a.d. sur $(\O,\A,\Pr)$.
    \begin{equation*}
        (X = Y) \nt{ presque sûr} \ra X \sim Y.
    \end{equation*}
    \tcblower
    Vérifions que $(X=Y)$ est un évènement.
    \begin{equation*}
        (X = Y) = \bigcup_{y \in Y(\O)}(X=y\cap Y=y) \in \A.
    \end{equation*}
    Vérifions que $X$ et $Y$ ont la même loi. Soit $x\in\R$.
    \begin{equation*}
        (X=x) = (X=x \cap X = Y) \sqcup (X=x\cap X\neq Y)
    \end{equation*}
    Or $\Pr(X=Y)=1-\Pr(X=Y)=0$ i.e. $X\neq Y$ est négligeable et $(X=x\cap X\neq Y)$ aussi.\\
    Donc $\Pr(X=x)=\Pr(X=x\cap X=Y) + \cancel{\Pr(X=x\cap X\neq Y)} = \Pr(Y=x\cap X=Y)$ donc $\Pr(X=x)=\Pr(Y=x)$.
\end{prop}

\begin{exercice}{}{}
    Soit $X$ une v.a.d. sur $(\O,\A,\Pr)$ à valeurs dans $\N$. Montrer que $\lim_{\infty}\Pr(X\geq n)=0$.
    \tcblower
    Déjà, posons la suite $(u_n)$ telle que $\forall n \in \N, ~ u_n = \Pr(X\geq n)$, décroissante puisque $(X\geq n+1)\subset(X\geq n)$.\\
    De plus, elle est minorée par zéro, elle converge donc par TLM. Notons $l$ sa limite.\\
    $\hra$ $(X\geq n) = \bigcup_{k\geq n}(X=k)$ donc $\Pr(X\geq n)=\sum_{k=n}^{\infty}\Pr(X=k)\to0$ en tant que reste de $\sum_{k=0}^\infty\Pr(X=k)=1$.\\
    $\hra$ Par continuité monotone, puisque $(u_n)$ est décroissante : $\lim\Pr(X\geq n)=\Pr(\bigcap_{n\in\N}(X\geq n))=0$.
\end{exercice}

\bf{Remarque.} Pour $(\O,\A,\Pr)$ un espace probabilisé et $X$ une v.a.d. sur cet espace, la loi de $X$ est bien définie ssi $\forall x \in X(\O), ~ \Pr(X=x)\geq0$ et $\sum_{x\in X(\O)}\Pr(X=x)=1$.

\begin{exercice}{}{}
    Soit $X$ une v.a.d. sur $(\O,\A,\Pr)$ à valeurs dans $\N$. Déterminer la loi de $X$ dans les cas suivants :
    \begin{enumerate}
        \item $\forall n \in \N^*, ~ \Pr(X=n)=\frac{2}{n}\Pr(X=n-1)$.
        \item $\exists k \in ~ ]0,1[, ~ \forall n \in \N, ~ \Pr(X=n)=k\Pr(X\geq n)$.
    \end{enumerate}
    \tcblower
    \boxed{1.} Posons $(u_n)$ la suite telle que $\forall n \in \N, ~ u_n = \Pr(X=n)$, alors $\forall n \in \N^*, ~ u_n = \frac{2}{n}u_{n-1}$.\\
    Alors $\forall n \in \N, ~ u_n = \frac{2^n}{n!}u_0$ (récurrence facile); et $\sum_{n=0}^\infty u_n = \Pr(X\in\N)=1$ donc :
    \begin{equation*}
        \sum_{n=0}^{\infty}\frac{2^n}{n!}u_0 = 1 \ra u_0 = e^{-2}.
    \end{equation*}
    Finalement, $\forall n \in \N, ~ \Pr(X=n)=\frac{2^n}{n!}e^{-2}$.\\
    \boxed{2.} Posons $(u_n)$ la suite telle que $\forall n \in \N, ~ u_n=\Pr(X=n)$ et $(v_n)$ telle que $\forall n \in \N, ~ v_n = \Pr(X\geq n)$.\\
    Alors $\forall n \in \N, ~ u_n = kv_n$ et $(X\geq n) = (X = n)\cup(X\geq n+1)$ donc $v_n=u_n + v_{n+1}=kv_n+v_{n+1}$.\\
    Ainsi, $v_{n+1}=(1-k)v_n$ donc $\forall n \in \N, ~ v_n=(1-k)^nv_0$, or $v_0=\Pr(X\geq0)=1$ donc $v_n=(1-k)^n$.\\
    Puis $u_n = kv_n = k(1-k)^n$. C'est une loi de probabilité : $\forall n \in \N, ~ u_k \geq 0$ et $\sum_{n=0}^\infty k(1-k)^n\to\frac{k}{1-(1-k)}=1$.
\end{exercice}

\begin{prop}{$f(X)$}{}
    Soit $X$ une v.a.d. sur $(\O,\A,\Pr)$ et $f:X(\O)\to\R$. On définit $f(X):w\mapsto f(X(\w))$ de $\O$ dans $\R$.\\
    Alors $f(X)$ est encore une v.a.d. sur $(\O,\A,\Pr)$.
    \tcblower
    Soit $X$ une v.a.d. sur $(\O,\A,\Pr)$. Déjà, $f(X(\O))$ est au plus dénombrable en tant qu'image par une application d'un ensemble au plus dénombrable puis $\forall y \in \R, ~ (f(X)=y) = (X \in f^{-1}(\{y\}))\in\A$.
\end{prop}

\vspace*{-0.8cm}

\section{Lois usuelles.}

\begin{defi}{Loi de Bernoulli}{}
    Soit $X$ une v.a.d. sur $(\O,\A,\Pr)$. On dit que $X$ suit une loi de Bernoulli de paramètre $p\in[0,1]$ ssi :\\[-0.4cm]
    \begin{equation*}
        X(\O) = \{0,1\}, ~ \Pr(X=1)=p \et \Pr(X=0) = 1-p.
    \end{equation*}
    On note alors $X\sim\B(p)$.
\end{defi}

\vspace*{-0.8cm}

\bf{Remarque.} Si $X\sim\B(p)$ alors $X^2=X$ puis $\forall n \in \N^*, ~ X^n = X$.

\begin{prop}{}{}
    Soit $X$ une v.a.d. sur $(\O,\A,\Pr)$.
    \begin{equation*}
        X \sim \B(p) ~ \nt{avec} ~ p\in[0,1] \iff X=\1_a ~ \nt{avec} ~ A\in\A.
    \end{equation*}
    \tcblower
    \boxed{\la} $(\1_A=1)=A$ et $(\1_A = 0) = \ov{A}$ donc $\1_A\sim\B(p)$ avec $p=\Pr(A)\in[0,1]$.\\
    \boxed{\ra} Soit $A=(X=1)\in\A$, alors $(X=1)\ra \1_A=1$, sinon $(X=0)\ra\1_A=0$.
\end{prop}

\vspace*{-0.8cm}

\begin{defi}{Loi binomiale.}{}
    Soit $X$ une v.a.d. sur $(\O,\A,\Pr)$. On dit que $X$ suit une loi binomiale de paramètres $p\in[0,1]$ et $n\in\N^*$ ssi :\\[-0.3cm]
    \begin{equation*}
        X(\O) = \lb0,n\rb, \et \forall k \in \lb0,n\rb, ~ \Pr(X=k) = \binom{n}{k}p^k(1-p)^{n-k}.
    \end{equation*}
    On note alors $X\sim \B(n,p)$.
    \tcblower
    Vérifions que la loi de $X$ est bien définie :
    \begin{equation*}
        \forall k \in \lb0,n\rb, ~ \binom{n}{k}p^k(1-p)^{n-k} \geq 0 \et \sum_{k=0}^n\binom{n}{k}p^k(1-p)^{n-k}=(p+(1-p))^n = 1.
    \end{equation*}
\end{defi}

\vspace*{-0.8cm}

\begin{defi}{Loi géométrique.}{}
    Soit $X$ une v.a.d. sur $(\O,\A,\Pr)$. On dit que $X$ suit une loi géométrique de paramètre $p\in~]0,1[$ ssi :\\[-0.3cm]
    \begin{equation*}
        X(\O) = \N^* \et \forall k \in \N^*, ~ \Pr(X=k) = p(1-p)^{k-1}.
    \end{equation*}
    On note alors $X\sim\cursive{G}(p)$.
    \tcblower
    La loi de $X$ est bien définie : $\forall k \in \N^*, ~ p(1-p)^{k-1} \geq 0 \et \sum_{k=1}^\infty p(1-p)^{k-1} = \sum_{k=0}^\infty(1-p)^k = \frac{p}{1-(1-p)}=1$.
\end{defi}

\vspace*{-0.8cm}

\begin{inter}{}{}
    On note $A$ un évènement succès sur $(\O,\A,\Pr)$, on répète des tirages indépendants jusqu'au succès.\\
    On note $X$ le nombre de tirages pour accéder au premier succès et $B_i$ : << succès au $i^{\nt{ème}}$ tirage >>.\\
    On a $X$ à valeurs dans $\N^*$ car il faut jouer au moins une fois pour réussir puis $(X=k)=\ov{B_1}\cap...\cap\ov{B_{k-1}}\cap B_k$.\\
    Les tirages sont indépendants donc la famille $(B_n)_{n\in\N^*}$ est mutuellement indépendante.
    \begin{equation*}
        \Pr(X=k)=\Pr(\ov{B_1}) \times ... \times \Pr(\ov{B_{k-1}}) \times \Pr(B_k) = \Pr(\ov{A})\times...\times \Pr(\ov{\ov{A}})\times \Pr(A) = (1-p)^{k-1}p.
    \end{equation*}
\end{inter}

\vspace*{-0.8cm}

\begin{prop}{}{}
    Soit $X$ une v.a.d. sur $(\O,\A,\Pr)$ à valeurs dans $\N^*$ et $p\in~]0,1[$.
    \begin{equation*}
        X \sim \cursive{G}(p) \iff \forall n \in \N, ~ \Pr(X > n) = (1-p)^n 
    \end{equation*}
    \tcblower
    \boxed{\ra} $(X>n) = \bigcup_{k\geq n+1}(X=k)$ indépendants 2-à-2 donc :
    \begin{equation*}
        \Pr(X>n) = \sum_{k=n+1}^{\infty} \Pr(X=k) = \sum_{k=n+1}^\infty p(1-p)^{k-1} = p(1-p)^n \sum_{k=0}^\infty(1-p)^k = \frac{p(1-p)^n}{p}=(1-p)^n.\\
    \end{equation*}
    \boxed{\la} $\forall n \in \N^*, ~ (X>n-1)=(X=n)\cup(X>n)$ puis $\Pr(X>n-1)=\Pr(X=n)+\Pr(X>n)$.
    \begin{equation*}
        \forall n \in \N^*, ~ \Pr(X=n)=\Pr(X>n-1) - \Pr(X>n) = (1-p)^{n-1} - (1-p)^n = p(1-p)^{n-1}
    \end{equation*}
    On a bien $X\sim\cursive{G}(p)$.
\end{prop}

\vspace*{-0.8cm}

\begin{prop}{Lois sans mémoire.}{}
    Soit $X$ une v.a.d. sur $(\O,\A,\Pr)$ à valeurs dans $\N^*$ et $p\in~]0,1[$.
    \begin{equation*}
        X \sim \cursive{G}(p) \iff \left[ \forall k,n \in \N, ~ \Pr(X>k+n \mid X>n) = \Pr(X>k)\right]
    \end{equation*}
    Alors $p=\Pr(X=1)$.
    \tcblower
    \boxed{\ra} $\forall k,n \in \N, ~ \Pr(X>k+n\mid X>n) = \frac{\Pr(X>k+n ~ \cap ~ X > n)}{\Pr(X>n)}=\frac{\Pr(X > k + n)}{\Pr(X > n)} = \frac{(1-p)^{k+n}}{(1-p)^n}=(1-p)^k=\Pr(X>k)$.\\
    \boxed{\la} $\Pr(X>n+1\mid X>n)=\frac{\Pr(X>n+1)}{\Pr(X>n)}=\Pr(X>1)$ par les calculs précédents.\\
    Donc $\Pr(X>n+1)=\Pr(X>1)\Pr(X>n)=(1-p)\Pr(X>n)$ puis $\forall n \in \N, ~ \Pr(X>n)=(1-p)^n\cancel{\Pr(X>0)}$.
\end{prop}

\vspace*{-0.8cm}

\begin{defi}{}{}
    Soit $X$ une v.a.d. sur $(\O,\A,\Pr)$ et $\l\in \R_+^*$. On dit que $X$ suit une loi de Poisson de paramètre $\l$ ssi\\[-0.3cm]
    \begin{equation*}
        X(\O) = \N \et \forall n \in \N, ~ \Pr(X=n) = e^{-\l}\frac{\l^n}{n!}.
    \end{equation*}
    On note alors $X\sim\cursive{P}(\l)$.
    \tcblower
    La loi de $X$ est bien définie :
    \begin{equation*}
        \forall n \in \N, ~ e^{-\l}\frac{\l^n}{n!} \geq 0 \et \sum_{n=0}^\infty e^{-\l}\frac{\l^n}{n!}=e^{-\l}e^\l = 1.
    \end{equation*}
\end{defi}

\begin{prop}{}{}
    Soit $(X_n)_n$ une suite de v.a.d. sur $(\O,\A,\Pr)$ telle que $\forall n \in \N, ~ X_n \sim \B(n,p_n)$ et $np_n \xrightarrow[n\to+\infty]{}\l>0$.\\
    Alors $\forall k \in \N, ~ \Pr(X_n = k) \xrightarrow[n\to+\infty]{}e^{-\l}\frac{\l^k}{k!}$.
    \tcblower
    Soit $(n,k)\in\N^2$, alors $\Pr(X_n=k)=\binom{n}{k}p_n^k(1-p_n)^{n-k}\sim_{\infty} \frac{n^k}{k!}\cdot\left( \frac{\l}{n} \right)^k\cdot e^{-\l} = e^{-\l}\frac{\l^k}{k!}$.\\
    On a $(1-p_n)^{n-k}=\frac{(1-p_n)^n}{(1-p_n)^k}\sim_\infty(1-p_n)^n=e^{n\ln(1-p_n)}$ puis $n\ln(1-p_n)\sim-np_n\to-\l$ puis cont. de $\exp$ en -$\l$.
\end{prop}

\section{Vecteurs aléatoires, couples de variables aléatoires.}


\begin{defi}{Vecteur aléatoire.}{}
    Soit $(X_n)$ une suite de v.a.d. sur $(\O,\A,\Pr)$. Un vecteur aléatoire discret $V$ est une application de $\O$ dans $\R^n$ tel que $V=(X_1,...,X_n)$ où $\forall i \in \lb1,n\rb, ~ X_i$ est une v.a.d. sur $(\O,\A,\Pr)$.
\end{defi}

\begin{defi}{Loi conjointe.}{}
    Soient $X,Y$ deux v.a.d. et $V=(X,Y)$. La loi de $V$, ou loi conjointe de $(X,Y)$ est la donnée de :\\
    $\hra$ $V(\O)$ l'ensemble des valeurs possibles de $V$, et $\Pr((X=x)\cap(Y=y))$ pour tout couple $(x,y)\in X(\O)\times Y(\O)$.
\end{defi}

\begin{exercice}{}{}
    Soient $X,Y$ deux v.a.d. sur $(\O,\A,\Pr)$ données par leur loi conjointe.\\
    Déterminer $\Pr(X=Y)$, $\Pr(X\geq Y)$, $\Pr(X+Y=a)$.
    \tcblower
    \begin{equation*}
        (X=Y) = \bigcup_{x\in X(\O)}(X=x\cap Y=x) \ra \Pr(X=Y) = \sum_{x\in X(\O)}\Pr(X=x \cap Y=x).
    \end{equation*} 
    \begin{equation*}
        (X\geq Y) = \bigcup_{x\geq y}(X=x\cap Y = y) \ra \Pr(X\geq Y) = \sum_{(x,y)\in V(\O)}\1_{x\geq y}\Pr(X=x\cap Y=y).
    \end{equation*}
    \begin{equation*}
        \Pr(X+Y=a) = \sum_{x\in X(\O)} \Pr(X=x \cap Y=a-x).
    \end{equation*}
\end{exercice}

\begin{defi}{Loi marginale.}{}
    Ce sont les lois de $X,Y$, déduites de celles du couple $(X,Y)$ par :
    \begin{equation*}
        \Pr(X=x) = \sum_{y\in Y(\O)}(\Pr(X=x))
    \end{equation*}
\end{defi}

\begin{exercice}{}{}
    Soient $X,Y$ v.a.d. sur $(\O,\A,\Pr)$ à valeurs dans $\N$.\\
    On suppose que $X\sim\cursive{P}(\l)$ où $\l>0$ et que $(Y\mid X=n)\sim \B(n,p)$.\\
    Déterminer la loi conjointe de $(X,Y)$ et reconnaître la loi marginale de $Y$.
    \tcblower
    Par hypothèse, $\forall n \in \N, ~ \Pr(X=n)=e^{-\l}\frac{\l^n}{n!}$ et $\forall k \in \N, ~ \Pr(Y=k\mid X=n)=\binom{n}{k}p^k(1-p)^{n-k}$ si $k\leq n$.\\
    Alors $\forall k \in \N, ~ \forall n \in \N, ~ \Pr(Y=k \cap X=n) = \Pr(Y=k\mid X=n)\Pr(X=n)=\1_{k\leq n}\binom{n}{k}p^k(1-p)^{n-k}e^{-\l}\frac{\l^n}{n!}$.\\
    Retrouvons la loi de $Y$ :
    \begin{align*}
        \forall k \in \N, ~ \Pr(Y=k) &= \sum_{n=0}^{\infty}\Pr(Y=k \cap X=n) = \sum_{n=k}^\infty \binom{n}{k}p^k(1-p)^{n-k} e^{-\l}\frac{\l^n}{n!}\\
        &=e^{-\l}\frac{p^k\l^k}{k!}\sum_{n=k}^{\infty}\frac{(1-p)^{n-k}\l^{n-k}}{(n-k)!} = e^{-\l}\frac{(p\l)^k}{k!}\sum_{n=0}^{\infty}\frac{(\l(1-p))^n}{n!}\\
        &=e^{-\l p}\frac{(\l p)^k}{k!} \ra Y \sim \cursive{P}(\l p).
    \end{align*}
\end{exercice}

\begin{prop}{}{}
    \begin{equation*}
        X \sim Y ~ \cancel{\ra} ~ X+X \sim X+Y
    \end{equation*}
    Exemple : $X\sim\B(1/2)$ et $Y=1-X\sim\B(1/2)$.
\end{prop}

\section{Indépendance de variables aléatoires discrètes.}

\begin{defi}{}{}
    Soient $X,Y$ deux v.a.d. sur $(\O,\A,\Pr)$, elles sont indépendantes ($X \ind Y$) ssi :
    \begin{equation*}
        \forall (x,y) \in X(\O) \times Y(\O), ~ \Pr(X = x \cap Y = y) = \Pr(X=x)\Pr(Y=y).
    \end{equation*}
\end{defi}

\begin{thm}{}{}
    Soient $X,Y$ deux v.a.d. sur $(\O,\A,\Pr)$. Elles sont indépendantes ssi :
    \begin{equation*}
        \forall A \subset X(\O), ~ \forall B \subset Y(\O), ~ \Pr(X\in A \cap Y \in B) = \Pr(X\in A)\Pr(Y \in B).
    \end{equation*}
    \tcblower
    \boxed{\la} Soit $x\in X(\O)$ et $y\in Y(\O)$. On pose $A=\{x\}\subset X(\O)$ et $B=\{y\}\subset Y(\O)$, puis $(X \subset \{x\}) = (X=x)$.\\
    \boxed{\ra} Soit $A\subset X(\O)$ et $B\subset Y(\O)$. On a $(X\in A \cap Y \in A) = \bigcup_{\substack{x\in A\\y\in B}}(X=x\cap Y=y)$. Par $\s$-additivité :
    \begin{equation*}
        \Pr(X\in A \cap Y \in B) = \sum_{\substack{x\in A\\y\in B}}\Pr(X=x\cap Y=y)=\sum_{\substack{x\in A\\ y \in B}}\Pr(X=x)\Pr(Y=y)
    \end{equation*}
    Puis par Fubini positif : $\Pr(X\in A \cap Y \in B) = \sum_{x\in A}\Pr(X=x)\sum_{y\in B}\Pr(Y=y) = \Pr(X\in A)\Pr(Y\in B)$.
\end{thm}

\begin{defi}{}{}
    Soient $X_1,...,X_n$ des v.a.d. sur $(\O,\A,\Pr)$. Il y a équivalence :
    \begin{enumerate}
        \item Elles sont indépendantes.
        \item $\forall (x_1,...,x_n) \in (X_1(\O)\times...\times X_n(\O)), ~ \Pr(X_1=x_1 \cap ... \cap X_n=x_n)=\Pr(X_1=x_1)...\Pr(X_n=x_n)$.
        \item $\forall (A_1,...,A_n) \in (\P(X_1(\O)),...,\P(X_n(\O))), ~ \Pr(X_1\in A_1\cap...\cap X_n\in A_n)=\Pr(X_1\in A_1)...\Pr(X_n\in A_n)$.
    \end{enumerate}
    Une famille quelconque est indépendante ssi toute sous-famille finie est indépendante. 
\end{defi}

\begin{prop}{}{}
    Soient $X,Y$ deux v.a.d. sur $(\O,\A,\Pr)$, $f:X(\O)\to\R$ et $g:Y(\O)\to\R$.\\
    Alors $X\ind Y \ra f(X) \ind g(Y)$.
    \tcblower
    On a $\forall A,B \subset \R, ~ \Pr(f(x)\in A \cap g(y)\in B)=\Pr(X\in f^{-1}(A))\Pr(Y\in g^{-1}(B))=\Pr(f(X)\in A)\Pr(g(Y)\in B)$.
\end{prop}

\begin{lemme}{Coalitions.}{}
    Soit $(X_i)_{i\in I}$ une famille de v.a.d. indépendantes sur $(\O,\A,\Pr)$ et $J,K\subset I$ \bf{disjointes}.\\
    Alors toute fonction de $(X_i)_{i\in J}$ est indépendante de toute fonction de $(X_i)_{i\in K}$.
\end{lemme}

\begin{prop}{Somme de variables.}{}
    Soient $X_1,...,X_n$ des v.a.d. indépendantes sur $(\O,\A,\Pr)$.
    \begin{enumerate}
        \item Si $\forall i \in \lb1,n\rb, ~ X_i \sim \B(p)$, alors $Z=\sum_{i=1}^n X_i \sim \B(n,p)$.
        \item Si $\forall i \in \lb1,n\rb, ~ X_i \sim \cursive{P}(\l_i)$, alors $Z=\sum_{i=1}^nX_i \sim \cursive{P}\left( \sum_{i=1}^n \l_i \right)$.
    \end{enumerate}
    \tcblower
    \boxed{1.} MP2I Espaces probabilisés, proposition 54.\\
    \boxed{2.} Par récurrence sur $n$. Vrai pour $n=1$. Pour $n=2$ :\\
    Soient $X,Y$ telles que $X\sim\cursive{P}(\l)$ et $Y\sim\cursive{P}(\mu)$, alors $X\ind Y$ et $X+Y$ est à valeurs dans $\N$.\\
    Alors $(X+Y = n) = \bigcup_{k=0}^n(X = k \cap Y = n - k)$ donc : \begin{equation*} \Pr(X+Y=n)=\sum_{k=0}^n\Pr(X=k)\Pr(Y=n-k)=\frac{e^{-(\l+\mu)}}{n!}\sum_{k=0}^n\binom{n}{k}\l^k\mu^{n-k} = e^{-\l + \mu} \frac{(\l+\mu)^n}{n!} \end{equation*}
    On reconnaît une loi de Poisson.\\
    Hérédité : Soient $X_1,...,X_{n+1}$ indépendantes alors par le lemme, $X_1+...+X_n$ et $X_1+...+X_{n+1}$ sont indépendantes.\\
    D'après la proposition au rang $n$ : $X_1+...+X_n \sim \Pr(\l_1+...+\l_n)$ puis par 2 : $(X_1+...+X_n)+X_{n+1}\sim\cursive{P}((\l_1+...+\l_n) + \l_{n+1})$.
\end{prop}

\pagebreak

\begin{exercice}{}{}
    \begin{enumerate}
        \item Soient $X,Y$ deux v.a.d. sur $(\O,\A,\Pr)$ avec $X$ constante presque sûrement. Montrer que $X\ind Y$.
        \item Soit $X$ v.a.d. sur $(\O,\A,\Pr)$. Montrer que $X\ind X$ ssi $X$ constante presque sûrement.
    \end{enumerate}
    \tcblower
    \boxed{1.} Par hypothèse, il existe $\l\in\R, ~ \Pr(X=\l)=1$, alors $\forall \mu \neq \l, ~ \Pr(X=\mu)=0$.\\
    Si $x\neq\l$ et $y\in Y(\O)$, $(X=x\cap Y=y) \subset (X=x)$, négligable : $\Pr(X=x\cap Y=y)=\Pr(X=x)\Pr(Y=y)=0$.\\
    Sinon, $x=\l$ et $\left((X=\l),(X\neq\l)\right)$ est un s.c.e. donc $\Pr(Y=y)=\Pr(Y=y\cap X=\l)+\cancel{\Pr(Y=y\cap X\neq \l)}.$\\
    Or $\Pr(X=\l)=1$ donc $\Pr(X=\l)\Pr(y=y)=\Pr(X=\l\cap Y=y)$.\\
    \boxed{2. \la} Cas particulier de la question précédente.\\
    \boxed{2. \ra} Par contraposée, on suppose qu'il existe $\l,\mu\in X(\O), ~ \l \neq \mu \et \Pr(X=\l)>0\et \Pr(X=\mu)>0$.\\
    Alors $\Pr(X=\l\cap X=\mu)=0$ mais $\Pr(X=\l)\Pr(X=\mu)>0$, donc $X$ n'est pas indépendante de $X$.
\end{exercice}

\vspace*{-0.9cm}

\begin{exercice}{}{}
    Soient $X,Y$ deux v.a.d. indépendantes sur $(\O,\A,\Pr)$ telles que $X\sim\B(p)$ et $Y\sim\B(q)$.\\
    Reconnaître la loi de $XY$.
    \tcblower
    On a $X(\O)=Y(\O)=\{0,1\}$, $\Pr(X=0)=1-p$, $\Pr(X=1)=p$, $\Pr(Y=0)=1-q$, $\Pr(Y=1)=q$.\\
    Donc $XY\sim\B(r)$ où $r=\Pr(XY=1)=\Pr(X=1\cap Y=1)=\Pr(X=1)\Pr(Y=1)=pq$.
\end{exercice}

\vspace*{-0.9cm}

\begin{exercice}{}{}
    Soient $A,B$ deux évènements sur $(\O,\A,\Pr)$. Montrer que $A\ind B \iff \1_A \ind \1_B$.
    \tcblower
    \boxed{\la} $\Pr(A\cap B)=\Pr((\1_A = 1) \cap (\1_B = 1))=\Pr(\1_A=1)\Pr(\1_B=1)=\Pr(A)\Pr(B)$.\\
    \boxed{\ra} Par hypothèse, $\Pr((\1_A = 1) \cap (\1_B = 1))=\Pr(\1_A=1)\Pr(\1_B=1)$.\\
    De plus, on a $A\ind B \iff \ov{A}\ind B \iff A \ind \ov{B} \iff \ov{A} \ind \ov{B}$.\\
    Donc $\Pr(\1_A = 0 \cap \1_B = 1) = \Pr(\ov{A}\cap B) = \Pr(\ov{A})\Pr(B) = \Pr(\1_A = 0)\Pr(\1_B=1)$...
\end{exercice}

\vspace*{-0.9cm}

\begin{exercice}{}{}
    Soit $(X_n)_{n\in\N^*}$ une suite de v.a.d. indépendantes sur $(\O,\A,\Pr)$, de mêmes lois que $X$.\\
    On note $S_n = \ds\sum_{k=1}^n X_k$ pour tout $n\in\N^*$.
    \begin{enumerate}
        \item Soit $m\geq n$. Montrer que $S_n$ et $S_{n+m}-S_m$ sont indépendantes et de même loi.
        \item Montrer que $\forall a \in \R, ~ (S_n \geq na) \subset \bigcup_{k=1}^n (X_k\geq a)$.
        \item Montrer que $\forall a \in \R, ~ \Pr(X\geq a) = 0 \iff \left[ \forall n \in \N^*, ~ \Pr(S_n \geq na) = 0 \right]$.
    \end{enumerate}
    \tcblower
    \boxed{1.} On a $S_n = X_1 + ... + X_n$ et $S_{n+m} - S_m = X_{m+1} + ... + X_{n+m}$ avec $m+1>n$.\\
    Or les $(X_n)$ sont indépendantes donc d'après le lemme des coalitions, $S_n$ et $S_{n+m}-S_m$ sont indépendantes.\\
    Soit $x\in\R$. $(S_n = x) = (X_1 + ... + X_n = x) = \bigcup_{x_1+...+x_n=x}(X_1=x_1\cap...\cap X_n=x_n)$.\\
    Alors $\Pr(S_n=x)=\sum_{x_1+...+x_n=x}\Pr(X_1=x_1)...\Pr(X_n=x_n)$ par indépendance.\\
    Puis $\forall i \in \lb1,n\rb, ~ \Pr(X_i=x_i)=\Pr(X_{m+i}=x_i)$ car mêmes lois donc $\Pr(S_n=x)=\Pr(S_{n+m}-S_m=x)$.\\
    \boxed{2.} Soit $a\in\R$, $\forall k \in \lb1,n\rb, ~ X_k < a \ra S_n < na$ donc par contraposée, $S_n\geq na \ra \exists k \in \lb1,n\rb, ~ X_k \geq a$.\\
    Donc $(S_n\geq na) \subset \bigcup_{k=1}^n (X_k \geq a)$.\\
    \boxed{3.} Soit $a\in\R$.\\
    \boxed{\la} Pour $n=1$, $\Pr(S_1 \geq a) = 0$ donc $\Pr(X\geq a)=0$.\\
    \boxed{\ra} Soit $n\in\N^*$, par 2) : $\Pr(S_n\geq na) \leq \Pr(\bigcup_{k=1}^n(X_k\geq a)) \leq \sum_{k=1}^n\Pr(X_k \geq a)=0$.\\
    Par encadrement, $\Pr(S_n \geq na) = 0$.
\end{exercice}

\begin{exercice}{}{}
    Soient $X\ind Y$ sur $(\O,\A,\Pr)$ telles que $X\sim\cursive{G}(p)$ et $Y\sim\cursive{G}(q)$.\\
    Quelle est la probabilité que $A=\begin{pmatrix}X&1\\0&Y\end{pmatrix}$ puis $B=\begin{pmatrix}X&-Y\\Y&-X\end{pmatrix}$ soit diagonalisable ?
    \tcblower
    \fbox{A.} Vérifions que $(A ~ \nt{diagonalisable})$ est un évènement.\\
    Pour $w\in\O$, $A$ est diagonalisable ssi $X(\w)\neq Y(\w)$ car alors $\X_A$ est SRS, et puisque si $X(\w)=Y(\w)$, on a $\Sp(A)=\{X(\w)\}$ donc $A$ diagonalisable impliquerait $A=X(\w)I_2$, absurde.\\
    Donc $(A~\nt{diagonalisable})=(X\neq Y)\in\A$. On a $(X=Y)=\bigcup_{x\in X(\O)}(X=x\cap Y=x)$.\\
    Donc $\Pr(X=Y)=\sum_{x\in X(\O)}\Pr(X=x)\Pr(Y=x)=\sum_{n=1}^{+\infty}p(1-p)^{n-1}q(1-q)^{n-1} = \frac{pq}{1-(1-p)(1-q)}$.\\
    Donc $\Pr(A~\nt{diagonalisable})=1-\Pr(X=Y)=\frac{p+q-2pq}{p+q-pq}$.\\
    \fbox{B.} Soit $\w\in\O$. $B$ est diagonalisable ssi $X(\w)>Y(\w)$ car alors $\X_B$ est SRS, et puisque si $X(\w)<Y(\w)$, on a $\X_B = X^2 - (X(\w)^2 - Y(\w)^2)$ non scindé, et si $X(\w)=Y(\w)$, alors $\X_B = X^2$ donc $B$ est nilpotente donc $B=0$, absurde car la loi géométrique est à valeurs strict. positives.\\
    Finalement, $(B~\nt{diagonalisable})=(X>Y)$ et $(X>Y)=\bigcup_{y\in Y(\O)}(X>y\cap Y=y)$.\\
    Donc $\Pr(X>Y)=\sum_{\in Y(\O)}\Pr(X>y)\Pr(Y=y)=\sum_{n=1}^{+\infty}(1-p)^nq(1-q)^{n-1}=\frac{q(1-p)}{1-(1-p)(1-q)}$. 
\end{exercice}

\vspace*{-0.8cm}

\begin{exercice}{}{}
    Soient $X\ind Y$ sur $(\O,\A,\Pr)$ telles que $X\sim\cursive{P}(\l)$ et $Y\sim\cursive{P}(\mu)$.\\
    Reconnaître la loi de $X$ sachant que $X+Y=n$.
    \tcblower
    Il s'agit de calculer $\Pr(X=k\mid X+Y=n)$ pour tout $k\in\N^*$. Soit $k\in\N^*$.
    \begin{equation*}
        \Pr(X=k \mid X+Y=n) = \frac{\Pr(X=k \cap X+Y=n)}{\Pr(X+Y=n)} = \frac{\Pr(X=k)\Pr(Y=n-k)}{\Pr(X+Y=n)}.
    \end{equation*}
    Si $k>n$, alors $\Pr(X=k\cap X+Y=n)=(X=k\cap Y=n-k)=\0$.\\
    Sinon, on sait que $X+Y\sim \cursive{P}(\l + \mu)$, donc \begin{equation*}\Pr(X=k\mid X+Y=n)=\frac{e^{-\l}\l^k}{k!}\frac{e^{-\mu}\mu^{n-k}}{(n-k)!}\frac{n!}{e^{-(\l+\mu)}(\l+\mu)^n}=\binom{n}{k}\frac{\l^k\mu^{n-k}}{(\l+\mu)^n}\end{equation*}
    La loi conditionnelle est une loi binomiale de paramètre $n,p$ avec $p=\frac{\l}{\l+\mu}\in[0,1]$.
\end{exercice}

\vspace*{-0.8cm}

\begin{exercice}{}{}
    Soient $X\ind Y$ sur $(\O,\A,\Pr)$ telles que $X\sim\cursive{G}(p)$ et $Y\sim\cursive{G}(q)$.\\
    Déterminer la loi de $Z=\min(X,Y)$ et la reconnaître.
    \tcblower
    Vérifions que $Z$ est une v.a.d. sur $(\O,\A,\Pr)$ et $Z(\O)\in\N$ car $X,Y$ le sont.\\
    $\hra$ Pour tout $n\in\N, ~ (Z\geq n) = (X\geq n \cap Y \geq n)$, c'est un évènement.\\
    Puis $(Z\geq n)=(Z = n) \cup (Z \geq n+1)$ donc $(Z=n)$ est un évènement et $Z$ est bien une v.a.d.\\
    Alors $\Pr(Z = n) = \Pr(Z \geq n) - \Pr(Z \geq n+1) = \Pr(X \geq n)\Pr(Y \geq n) - \Pr(X \geq n+1)\Pr(Y\geq n+1)$ par ind.\\
    Donc $\Pr(Z=n)=(1-p)^{n-1}(1-q)^{n-1} - (1-p)^n(1-q)^n = \left[ 1 - (1-p)(1-q) \right]\left[(1-p)(1-q)\right]^{n-1}$.\\
    Donc $Z\sim\cursive{G}(1-(1-p)(1-q))$.
\end{exercice}

\end{document}
